Newcomers learn which ethical norms a workplace actually practices mostly by watching how colleagues handle ethically charged situations and what happens afterwards. This became a problem for many recent entrants to office-based work, including many members of Generation Z, whose entry period took place partly at a distance. This paper develops a conceptual model of ethical socialization under distributed entry, grounded in social learning theory and imprinting theory. It proposes that distributed entry reduces newcomers' clarity about enacted ethical norms by limiting their observational access to ethical conduct in the work unit, and that leaders who narrate the ethical reasoning behind decisions can offset part of this loss. It further proposes that distance weakens the transmission of a unit's enacted norms whatever their ethical quality, a loss where norms are ethical and a partial protection where misconduct is normal, and that these differences persist beyond entry and follow entry conditions rather than birth cohort. The paper reports no data; it suggests implications for onboarding and outlines designs for testing the five propositions in field and experimental settings.
